%%%%%%%%%%  scan idx 59  |  folio 46  %%%%%%%%%%
\SGAPrintPage{46}%
braic sheaf on $X$ whose sections on an open subset $V$ are the sections of $F$ on $U\cap V$. We readily find that $\overline F$ is \emph{quasi-coherent} (i.e. defined on every affine open subset $V$ -- hence here on all of $X$ -- by a module, of finite type or otherwise, over the coordinate ring of the open subset in question). This is immediate when $U$ itself is affine, for $\overline F$ is then defined on the open subset $V$ by $\Gamma(U\cap V,F)$ considered as a module over $\Gamma(V,\underline{O}_X)$. In the general case, $U$ is a finite union of affine open subsets $U_i$, so $\overline F$ is the kernel of an evident homomorphism of quasi-coherent sheaves:
\[
 \sum_i \overline{F|U_i}\longrightarrow \sum_{i,j}\overline{F|U_i\cap U_j}
\]
and is therefore itself quasi-coherent. As $X$ is assumed affine, $\overline F$ is therefore generated by its sections. As $\overline F|U=F$ is coherent, there exist finitely many sections of $\overline F$ generating $\overline F|U$. These sections generate a coherent subsheaf of $\overline F$, whose restriction to $U$ is isomorphic to $F$.
\begin{flushright}q.e.d.\end{flushright}

Now assume $X$ quasi-projective and \emph{nonsingular}. Then, by the syzygy theorem, we have, with the notation of the corollary to theorem 2.3, $\cF_0(X)=\cF(X)$. In this case we therefore have:
\sgaStableTarget{sga6.e00.eq.2-13}\begin{equation}\tag{2.13}
 \underline{K}(\xi(X))=\underline{K}(\cF(X)).
\end{equation}
We can express this isomorphism as follows:

\sgaStableTarget{sga6.e00.cor.2}\textbf{Corollary 2 to theorem 2.3.} \emph{Let $X$ be a nonsingular quasi-projective variety. For every abelian group $A$, there is a one-to-one correspondence between additive functions $f$ of vector bundles on $X$, with values in $A$, and additive functions $g$ of coherent algebraic sheaves on $X$, with values in $A$. To every $f$ corresponds the unique $g$ such that we have:}
%%%%%%%%%%  scan idx 60  |  folio 47  %%%%%%%%%%
\SGAPrintPage{47}%
\sgaStableTarget{sga6.e00.eq.2-14}\begin{equation}\tag{2.14}
 f(E)=g(\underline{O}_X(E))
\end{equation}
\emph{for every vector bundle $E$ on $X$, denoting by $\underline{O}_X(E)$ the sheaf of germs of regular sections of $E$ on $X$.}

To calculate $g(F)$ for an arbitrary coherent algebraic sheaf, we consider a finite resolution of $F$ by sheaves of the form $\underline{O}_X(E_i)$, and we then have:
\sgaStableTarget{sga6.e00.eq.2-14bis}\begin{equation}\tag{2.14 bis}
 g(F)=\sum_{i\geq0}(-1)^if(E_i).
\end{equation}

In particular, when $X$ is connected, we considered in \S 1 an additive function $\Ctil(E)$ of the vector bundle $E$, with values in the group $\Atil=\Atil(X)$. We deduce from it an additive function of coherent sheaves, say $\Ctil(F)$, and, by passing to components, functions $F\longrightarrow C^i(F)\in A^i(X)$. The $C^i(F)$ are called the \emph{Chern classes of the coherent sheaf} $F$; they are theoretically calculated from the Chern classes of vector bundles by (2.14 bis) (where, however, we must pass to multiplicative notation in the second member). As for the component of $\Ctil(F)$ along $\Z$, it is the \emph{rank} of the sheaf $F$, which corresponds to the rank of a module over an integral domain: when $X$ is affine (a case to which we can always reduce), $F$ is defined by a module of finite type over the coordinate ring of $X$, and we take the rank of this module\dots.

\sgaStableTarget{sga6.e00.sec.3.occ3}\subsection*{\S 3. Functorial generalities on $\underline{K}(X)$}

When $X$ is an algebraic space, we put, for brevity, $\underline{K}(X)=\underline{K}(\cF(X))$, and, when $X$ is nonsingular and quasi-projective, we identify this group with $\underline{K}(\xi(X))$ by corollary 2 to theorem 2.3. As $\underline{K}(\xi(X))$ is
%%%%%%%%%%  scan idx 61  |  folio 48  %%%%%%%%%%
\SGAPrintPage{48}%
not
only an abelian group, but a $\lambda$-ring, it follows by transport of structure that $\underline{K}(X)$ is also a $\lambda$-ring, and is therefore equipped with a multiplicative structure and maps $\lambda^i$. It is important to define these structures directly, without passing to vector bundles; this will be done below.

If $E$ is a vector bundle on $X$, we denote by $\gamma(E)$ its class in $\underline{K}(X)$; if $F$ is a coherent algebraic sheaf on $X$, $\gamma(F)$ likewise denotes its class in $\underline{K}(X)$. If $Y$ is an irreducible subvariety of $X$, we put $\gamma(Y)=\gamma(\underline{O}_Y)$, where $\underline{O}_Y$ denotes the sheaf of local rings of $Y$, considered as a coherent algebraic sheaf on $X$. We define the symbol $\gamma(Y)$ by linearity when $Y$ is an arbitrary cycle on $X$. When there is a risk of confusion, we shall write $\gamma_X$ for $\gamma$.

Let $X$ be an arbitrary algebraic space. We introduce a \emph{filtration} on $\underline{K}(X)$, denoting by $\underline{K}(X)^i$ the subgroup of $\underline{K}(X)$ generated by the $\gamma(F)$ with dim.supp.\ $F\leq n-i$ (where $n$ is the dimension of $X$). From Serre's ``dévissage lemma'', in the form given in [3], follows the result:

\sgaStableTarget{sga6.e00.prop.2-6}\textbf{Proposition 2.6.} \emph{The group $\underline{K}(X)^i$ is generated by the $\gamma(Y)$, where $Y$ runs through the set of subvarieties of $X$ of dimension $\leq n-i$. If $X$ is irreducible, we have $G^0(\underline{K}(X))=\Z$.}

The latter isomorphism is given by the homomorphism from $\underline{K}(X)$ to $\Z$ defined by means of the notion of \emph{rank} of a coherent algebraic sheaf.

Now assume $X$ \emph{nonsingular}. Following the developments
%%%%%%%%%%  scan idx 62  |  folio 49  %%%%%%%%%%
\SGAPrintPage{49}%
at the beginning of \S 2, we
define a multiplication on $\underline{K}(X)$ by putting
\sgaStableTarget{sga6.e00.eq.2-15}\begin{equation}\tag{2.15}
 \gamma(F)\gamma(G)=\sum_{i\geq0}(-1)^i\gamma(\underline{\Tor}_i(F,G))
\end{equation}
(the $\underline{\Tor}_i$ being taken, of course, with respect to the local rings of $X$). The exact sequence of the $\underline{\Tor}_i$ shows that formula (2.15) does define a biadditive map from $\underline{K}(X)\times\underline{K}(X)$ to $\underline{K}(X)$; associativity follows from the spectral sequence of multiple $\Tor$. We also verify in this way that the product of several factors can be calculated by the formula:
\sgaStableTarget{sga6.e00.eq.2-16}\begin{equation}\tag{2.16}
 \gamma(F_1)....\gamma(F_p)=\sum_{i\geq0}(-1)^i\gamma(\underline{\Tor}_i(F_1,\ldots,F_p))
\end{equation}
where, in the second member, the $\underline{\Tor}_i$ are ``simultaneous $\underline{\Tor}$'', calculated by simultaneous projective resolution of the arguments $F_i$.

If, in (2.15), we assume $G$ \emph{locally free}, we find:
\sgaStableTarget{sga6.e00.eq.2-15bis}\begin{equation}\tag{2.15 bis}
 \gamma(F)\gamma(G)=\gamma(F\otimes G),
\end{equation}
which shows in particular that $\gamma(\underline{O}_X)$ is the unit element of $\underline{K}(X)$, and that the canonical homomorphism from $\underline{K}(\xi(X))$ to $\underline{K}(X)$ is a ring homomorphism.

Formula (2.15) is evidently suggested by Serre's ``$\Tor$ formula'' in intersection theory. This formula proves the second part of the following proposition, the first following immediately from formula (2.15) and a calculation of local $\underline{\Tor}_i$:

\sgaStableTarget{sga6.e00.prop.2-7}\textbf{Proposition 2.7.} \emph{Let $X$ be a nonsingular variety, and $Y$ and $Z$ two cycles in $X$. If $Y$ and $Z$ meet transversely everywhere, we have:}
\sgaStableTarget{sga6.e00.eq.2-17}\begin{equation}\tag{2.17}
 \gamma(Y)\gamma(Z)=\gamma(Y.Z).
\end{equation}
\emph{If $Y$ and $Z$ are homogeneous and have respective dimensions $n-i$ and $n-j$, and if they meet without an excess component, we then have:}
%%%%%%%%%%  scan idx 63  |  folio 50  %%%%%%%%%%
\SGAPrintPage{50}%
\sgaStableTarget{sga6.e00.eq.2-18}\begin{equation}\tag{2.18}
 \gamma(Y)\gamma(Z)=\gamma(Y.Z)\qquad\mathrm{mod.}\ \underline K(X)^{i+j+1}.
\end{equation}

We
must take care that equality does not hold in general in (2.18). We shall see in \S 4 that, when $X$ is \emph{nonsingular and quasi-projective, the ring structure of $\underline K(X)$ is compatible with its filtration and that the associated graded ring $G(\underline K(X))$ is identified with a quotient of $A(X)$ by a torsion subgroup (perhaps always reduced to $0$\SGAAuthorFootnote{*}{No; for a counterexample cf.\ this Seminar XIV 4.7.})}.

I know how to define the $\lambda^i(F)$ directly in terms of sheaves only when the field $k$ has characteristic $0$ (this is the reason that has prevented us from proving the Riemann--Roch theorem when $k$ is not of characteristic $0$). Let $F$ be a coherent algebraic sheaf on $X$ on which the group $\mathfrak S_p$ of permutations of $p$ letters operates; we then denote by $(F)^{\mathrm{alt}}$ the set of the $x\in F$ such that $s.x=\varepsilon(s)x$ for every $s\in\mathfrak S_p$ ($\varepsilon(s)$ is the \emph{signature} of $s$). This is evidently a coherent algebraic subsheaf of $F$. This said, we have:

\sgaStableTarget{sga6.e00.thm.2-8}\textbf{Theorem 2.8.} \emph{Let $X$ be a nonsingular quasi-projective variety over the field} $k$ of characteristic $0$. \emph{Let $F$ be a coherent algebraic sheaf on $X$; we then have:}
\sgaStableTarget{sga6.e00.eq.2-19}\begin{equation}\tag{2.19}
 \lambda^p(\gamma(F))=\sum_{i\geq0}(-1)^i\gamma\left((\underline{\Tor}_i(F,F,\ldots,F))^{\mathrm{alt}}\right)
\end{equation}

N.B. It is clear that the group $\mathfrak S_p$ acts on the sheaf $\underline{\Tor}_i(F,\ldots,F)$.

\emph{Principle of the proof}: we consider a finite resolution $\underline L$ of $F$ by locally free sheaves $L_i$, and use the following general lemma:

